% 横版扉页：沿用竖版扉页的元素、字体、字号与配色，仅按 A4 横向重新安排纵向位置。
% 由 layout-landscape.tex 末尾加载，沿用其字体与颜色；书名、作者、英文书名与配图路径由入口提供。
% 路径均相对于项目根目录；绝对定位依赖 TikZ 的 remember picture，至少编译两轮（根目录 run.sh 已保证）。

\definecolor{LandscapeTitlePageInk}{HTML}{1D1D1F}
\definecolor{LandscapeTitlePageSecondary}{HTML}{6E6E73}
\definecolor{LandscapeTitlePageHairline}{HTML}{D1D1D6}
% 唯一强调色与代码高亮的蓝灰一致。
\colorlet{LandscapeTitlePageAccent}{LandscapeCodeBlue}
\newfontfamily\LandscapeTitlePageMediumLatin{PingFang SC Medium}
\newCJKfontfamily\LandscapeTitlePageMediumCJK{PingFang SC Medium}
% 提交信息的眉题：小号大写，加宽字距。
\newfontfamily\LandscapeTitlePageEyebrow{PingFang SC Medium}[LetterSpace=14]
% 提交说明原样输出：此字体族不启用 TeX 连字，-- 等不会被替换。
\newfontfamily\LandscapeTitlePagePlainLatin{PingFang SC}

% 配图区域：等比缩放后居中；标题区与提交信息的位置不随图片比例变化。
% 页高只有 210 mm，区域由竖版的 130 × 110 mm 收紧为 110 × 70 mm。
\newlength{\LandscapeTitlePageImageWidth}
\newlength{\LandscapeTitlePageImageHeight}
\setlength{\LandscapeTitlePageImageWidth}{110mm}
\setlength{\LandscapeTitlePageImageHeight}{70mm}

\makeatletter
% 元信息：英文书名与配图路径；不写或参数为空时省略对应内容。
\newcommand*{\EnglishTitle}[1]{\gdef\LandscapeTitlePage@EnglishTitle{#1}}
\newcommand*{\TitlePageImage}[1]{\gdef\LandscapeTitlePage@Image{#1}}
\let\LandscapeTitlePage@EnglishTitle\@empty
\let\LandscapeTitlePage@Image\@empty
% PDF 文档属性同步入口的书名与作者；入口未设置时不写入。
\let\LandscapeTitlePage@DefaultTitle\@title
\let\LandscapeTitlePage@DefaultAuthor\@author
\AtBeginDocument{%
  \ifx\@title\LandscapeTitlePage@DefaultTitle\else\hypersetup{pdftitle={\@title}}\fi
  \ifx\@author\LandscapeTitlePage@DefaultAuthor\else\hypersetup{pdfauthor={\@author}}\fi}
\makeatother

% 提交信息文件由 scripts/commit-info.sh 经根目录 run.sh 写入编译输出目录，
% 四行依次为日期、短哈希、提交说明、clean 或 dirty。按原样逐行读取，
% 不经 TeX 转义；文件不存在或缺少日期、哈希时不显示提交信息。
\ExplSyntaxOn
\ior_new:N \g__landscape_titlepage_ior
\str_new:N \g__landscape_titlepage_date_str
\str_new:N \g__landscape_titlepage_hash_str
\str_new:N \g__landscape_titlepage_subject_str
\str_new:N \g__landscape_titlepage_state_str
\bool_new:N \g__landscape_titlepage_commit_bool
\cs_new_protected:Npn \__landscape_titlepage_get:N #1
  {
    \ior_if_eof:NTF \g__landscape_titlepage_ior
      { \str_gclear:N #1 }
      {
        \ior_str_get:NN \g__landscape_titlepage_ior \l_tmpa_str
        \str_gset_eq:NN #1 \l_tmpa_str
      }
  }
\NewDocumentCommand \LandscapeTitlePageReadCommit { m }
  {
    \bool_gset_false:N \g__landscape_titlepage_commit_bool
    \file_if_exist:nT {#1}
      {
        \ior_open:Nn \g__landscape_titlepage_ior {#1}
        \__landscape_titlepage_get:N \g__landscape_titlepage_date_str
        \__landscape_titlepage_get:N \g__landscape_titlepage_hash_str
        \__landscape_titlepage_get:N \g__landscape_titlepage_subject_str
        \__landscape_titlepage_get:N \g__landscape_titlepage_state_str
        \ior_close:N \g__landscape_titlepage_ior
        \bool_lazy_or:nnF
          { \str_if_empty_p:N \g__landscape_titlepage_date_str }
          { \str_if_empty_p:N \g__landscape_titlepage_hash_str }
          { \bool_gset_true:N \g__landscape_titlepage_commit_bool }
      }
  }
% 默认读取本次编译的 <jobname>.commit；样张可在 \maketitle 前改读其他文件。
\AtBeginDocument { \LandscapeTitlePageReadCommit { \c_sys_jobname_str .commit } }
\NewDocumentCommand \LandscapeTitlePageIfCommitTF { +m +m }
  { \bool_if:NTF \g__landscape_titlepage_commit_bool {#1} {#2} }
% 眉题式三行：眉题“LAST COMMIT”；日期与哈希，dirty 时附圆角细框标签；提交说明。
% 不用分隔符串联，依靠换行、留白及字体差异区分标签与数值。
\NewDocumentCommand \LandscapeTitlePageCommitText { }
  {
    { \LandscapeTitlePageEyebrow \fontsize { 6.5 } { 8 } \selectfont LAST ~ COMMIT } \\ [ 1.5pt ]
    { \ttfamily \str_use:N \g__landscape_titlepage_date_str }
    \hspace { 0.8em }
    { \ttfamily \color { LandscapeTitlePageAccent } \str_use:N \g__landscape_titlepage_hash_str }
    \str_if_eq:VnT \g__landscape_titlepage_state_str { dirty }
      { \hspace { 0.6em } \usebox \LandscapeTitlePageDirtyBox }
    \str_if_empty:NF \g__landscape_titlepage_subject_str
      { \\ \LandscapeTitlePagePlainLatin \str_use:N \g__landscape_titlepage_subject_str }
  }
\ExplSyntaxOff
% 未提交修改标签：0.4 pt 浅灰圆角细框，与代码卡片同一视觉语言。
% 在扉页的 tikzpicture 之外预先排入盒子，避免嵌套 TikZ。
\newsavebox{\LandscapeTitlePageDirtyBox}
\newcommand*{\LandscapeTitlePageMakeDirtyBox}{%
  \sbox{\LandscapeTitlePageDirtyBox}{\tikz[baseline=(tag.base)]
    \node[draw=LandscapeTitlePageHairline,line width=0.4pt,rounded corners=1.1mm,
      inner xsep=1.3mm,inner ysep=0.45mm,text=LandscapeTitlePageInk,
      font=\sffamily\fontsize{6.5}{8}\selectfont] (tag) {含未提交修改};}}

% 版面以纸张中心为原点，元素顺序与竖版相同：标题区在上，配图居中偏下，提交信息贴近底边。
% 标题区三行的间距与竖版相同；书名基线距页顶 48 mm，提交信息底边距纸边 17 mm（同竖版）。
% 横版正文为双栏；扉页内容全部绝对定位，不受分栏影响。
\makeatletter
\renewcommand{\maketitle}{%
  \clearpage
  \thispagestyle{empty}%
  \LandscapeTitlePageMakeDirtyBox
  \null
  \begin{tikzpicture}[remember picture,overlay]
    \begin{scope}[shift={(current page.center)}]
      \node[anchor=base,inner sep=0pt,text=LandscapeTitlePageInk,
        font=\sffamily\bfseries\fontsize{44}{52}\selectfont] at (0,57mm) {\@title};
      \ifx\LandscapeTitlePage@EnglishTitle\@empty\else
        \node[anchor=base,inner sep=0pt,text=LandscapeTitlePageSecondary,
          font=\sffamily\fontsize{15}{18}\selectfont] at (0,43.5mm)
          {\LandscapeTitlePage@EnglishTitle};
      \fi
      \node[anchor=base,inner sep=0pt,text=LandscapeTitlePageInk,
        font=\LandscapeTitlePageMediumLatin\LandscapeTitlePageMediumCJK\fontsize{12}{15}\selectfont]
        at (0,32mm) {\@author};
      \ifx\LandscapeTitlePage@Image\@empty\else
        \node[inner sep=0pt] at (0,-17mm) {\includegraphics[
          width=\LandscapeTitlePageImageWidth,height=\LandscapeTitlePageImageHeight,
          keepaspectratio]{\LandscapeTitlePage@Image}};
      \fi
      % 底边对齐：提交说明换行时向上增长，不进入下边距。
      % flush center 使用 fil 胶居中；align=center 的有限伸缩会拉开 xeCJK 字距。
      \LandscapeTitlePageIfCommitTF{%
        \node[anchor=south,inner sep=0pt,text width=150mm,align=flush center,
          text=LandscapeTitlePageSecondary,font=\sffamily\fontsize{8}{11}\selectfont]
          at (0,-88mm) {\LandscapeTitlePageCommitText};
      }{}%
    \end{scope}
  \end{tikzpicture}%
  \clearpage}
\makeatother
